\begin{table}[H]
\centering
\caption{Cod: finite-duration floors.}
\label{tab:si-cod-floors}
\small
\begin{tabular}{llrrrr}
\toprule
policy & class & $n=5$ & 10 & 15 & 20 \\
\midrule
BAU (5 kt) & worst & 1450.0 & 1803.7 & 2062.3 & 2275.0 \\
flat 60 kt & worst & 1656.7 & 2188.8 & 2657.5 & 3177.4 \\
S1 & worst & 1656.7 & 2188.8 & 2657.5 & 3177.4 \\
cascade & worst & 1656.7 & 2188.8 & 2657.5 & 3177.4 \\
flat 120 kt & worst & 1891.7 & 2671.8 & 3562.3 & 5311.8 \\
BAU (5 kt) & q05 & 1298.7 & 1540.7 & 1697.8 & 1806.9 \\
flat 60 kt & q05 & 1499.6 & 1893.2 & 2193.4 & 2456.2 \\
S1 & q05 & 1499.6 & 1893.2 & 2193.4 & 2456.2 \\
cascade & q05 & 1499.6 & 1893.2 & 2193.4 & 2456.2 \\
flat 120 kt & q05 & 1727.6 & 2329.0 & 2898.1 & 3619.7 \\
BAU (5 kt) & q10 & 884.6 & 884.6 & 884.6 & 884.6 \\
flat 60 kt & q10 & 884.6 & 884.6 & 884.6 & 884.6 \\
S1 & q10 & 884.6 & 884.6 & 884.6 & 884.6 \\
cascade & q10 & 884.6 & 884.6 & 884.6 & 884.6 \\
flat 120 kt & q10 & 979.9 & 1028.6 & 1053.9 & 1067.2 \\
\bottomrule
\end{tabular}
\par\smallskip\begin{minipage}{0.97\linewidth}\footnotesize Lower boundary of the set of states from which $n$ consecutive floor years can be absorbed and the nominal (zero-residual) map then holds the safe set. For $C<g(\theta)$ this coincides with the $n$-step kernel boundary $b_n$.\end{minipage}
\end{table}
